%% ma: L10-B12-0012
%% lop: 10
%% chuong: 5
%% bai: 12
%% dang: 10.12.4
%% loai: TL
%% muc: TH
%% dap_an: "-"
%% nguon: "chuyen-de-cac-so-dac-trung-cua-mau-so-lieu-khong-ghep-nhom-toan-10-chuong-trinh-moi.pdf (GV. Phan Nhật Linh) (P01), Bài 01 Bài tập tự luận 12"
%% kien_thuc: "Lớp 10 Bài 12 (số gần đúng, sai số tuyệt đối, độ chính xác, sai số tương đối, quy tròn số)"
%% trang_thai: chinh-thuc
%% ghi_chu: "Lời giải gốc dùng nhầm 45 kg thay cho 43 kg và kết luận sai (phép đo thứ nhất); đúng là phép đo của Cường. Tự duyệt theo yêu cầu giáo viên 22/10/2026."
\begin{ex}
Bạn Lan đo được cân nặng là $43\text{ kg}\pm0{,}2\text{ kg}$. Bạn Cường đo được cân nặng là $65\text{ kg}\pm0{,}3\text{ kg}$. Trong hai phép đo trên, phép đo nào chính xác hơn?
\loigiai{Phép đo của Lan: $\delta_1\le\dfrac{0{,}2}{43}\approx0{,}465\%$.

Phép đo của Cường: $\delta_2\le\dfrac{0{,}3}{65}\approx0{,}462\%$.

Vì $\delta_2<\delta_1$ nên phép đo của bạn Cường chính xác hơn.}
\end{ex}
